\begin{theorem}[Complete metric limit uniqueness orphan RegSeqRat route]
\label{thm:complete-metric-limit-uniqueness-orphan-regseqrat-route}
Every $\CompleteMetricLimitUniquenessUp$ uniqueness handoff consumed by a separated-limit reader first exposes separated $\MetricSpaceUp$ identity together with the $\StreamNameUp$ finite-window rows, the $\RegSeqRatUp$ readback row, and the terminal $\RealUp$ zero-distance seal. The route is a finite child of the displayed packet
$$
(A,L_{0},L_{1},M,Z,S,H,C,P,N)
$$
and exports no host equality of metric points, selected limit, quotient representative, countable-choice schedule, or ambient $\RealUp$ completeness principle.
\end{theorem}
\begin{proof}
Project the carrier of \autoref{def:complete-metric-limit-uniqueness-carrier}. The two complete-metric limit rows, metric witness, zero-distance row, separated classifier, transport, replay, provenance, and local name already give the public comparison packet.

Use the regular-rational and Real-facing child routes, including \autoref{thm:complete-metric-limit-uniqueness-regseqrat-forward-route} and \autoref{thm:complete-metric-limit-uniqueness-regseqrat-real-separated-dependency-route}. They force the $\StreamNameUp$ windows, $\RegSeqRatUp$ readback, dyadic tolerance data, and $\RealUp$ zero-distance seal to be displayed before the separated classifier is consumed.

The separated route then reads identity only through the displayed metric and zero-distance rows. Since the carrier has no coordinate for host point equality, selected limits, quotient representatives, countable choice, or ambient Real completeness, the handoff remains the stated finite child route.
\end{proof}
